\section{从流程设计到上下文设计}

本期最重要的产品方法论，是从 flow design 转向 context design。互联网产品过去 20 年的基本范式，是产品经理预设用户路径、输入和输出；设计师画出界面；工程师实现流程。AI Native 产品里，用户输入开放，模型输出开放，尤其 Agent 时代还会调用工具、生成文件、执行动作，流程无法被完全穷举。

\begin{figure}[H]
\centering
\includegraphics[width=0.96\textwidth]{flow-to-context-design.png}
\caption{从流程设计到上下文设计：有限自由度转向开放输入与开放输出。}
\end{figure}

\begin{knowledgebox}{读图：AI Native 的设计对象变了}
左侧流程设计面向按钮、页面、漏斗和固定输出；右侧上下文设计面向 prompt、context、工具、模型组合、约束和反馈。产品流程仍然重要，但它要围绕模型能力和上下文效果反推，而不是先画完整流程再期待模型配合。
\end{knowledgebox}

\subsection{为什么传统协作会失效}

前面的图把范式差异画出来了，这里落到团队协作。AI Native 不是产品经理一个人的认知变化，而是产品、设计、工程和模型团队的协作顺序变化。

传统团队协作是线性的：产品经理写需求，设计师画图，工程师施工。AI Native 产品里，最关键的中间部分可能是用户看不见的 context：给模型什么材料、怎样组织任务、调用哪些工具、如何评测输出、怎样让用户反馈进入下一轮。这些东西不先跑通，UI 流程可能越精致越错。

\begin{warningbox}{流程先行的风险}
如果先把页面、按钮和用户路径全部设计完，再去接模型，很容易发现模型真实输出和预设流程不匹配。AI Native 产品应先验证能力闭环，再设计外层流程。
\end{warningbox}

\subsection{本章小结}

流程设计到上下文设计，是互联网产品经理转向 AI 产品的关键认知鸿沟。产品经理不再只是用户路径设计师，而要成为模型能力、上下文、反馈和用户心智的组织者。

